% Curated cumulative working-reader wrapper through HI-OLP-0011.
% Source authority: OpenLogicProject/OpenLogic commit
% 9620cc73f9c8e0ad003c514a5d3748f29611c4c0.
% This wrapper adds only Hindi title/status/navigation matter; each translated
% section remains an independently hashed source-preserving subfile.

\documentclass[../../../include/open-logic-section]{subfiles}
\usepackage{flafter}
\usepackage{float}

\renewcommand*{\ifgitinfo}{\sffamily\footnotesize संशोधन~9620cc7}
\renewcommand*{\photocredits}{}
\renewcommand*{\contentsname}{विषय-सूची}
% The generated job name is intentionally descriptive on disk, but it is too
% long for the running footer.  Keep the linked work name and licence visible
% while using a compact display label that cannot collide with the page number.
\renewcommand*{\olshortlicensetext}[1]
  {\textit{\href{\OLTlinkname}{OpenLogic-hi}}
  \ifgitinfo\ 
  #1 द्वारा /
  \href{http://creativecommons.org/licenses/by/4.0/}{CC--BY}.}
\floatplacement{figure}{H}
\setlength{\emergencystretch}{2em}

\hypersetup{
  pdftitle={ओपन लॉजिक परियोजना: हिन्दी संचयी अनुवाद},
  pdfauthor={Interlanguage Hindi translation programme},
  pdfsubject={समुच्चयों और संबंधों पर ग्यारह अनूदित खंड},
  pdfkeywords={तर्कशास्त्र, समुच्चय, संबंध, आलेख, हिन्दी, Open Logic Project}
}

\begin{document}

\begingroup
\thispagestyle{empty}
\centering
\vspace*{0.11\textheight}
{\Huge\bfseries ओपन लॉजिक परियोजना\par}
\vspace{0.8cm}
{\LARGE हिन्दी संचयी अनुवाद\par}
\vspace{1.0cm}
{\Large समुच्चय और संबंध\par}
\vspace{0.4cm}
{\large ग्यारह आरम्भिक खंड\par}
\vfill
{\normalsize स्रोत: Open Logic Project, संशोधन
\texttt{9620cc73f9c8e0ad003c514a5d3748f29611c4c0}\par}
\vspace{0.35cm}
{\normalsize स्थायी हिन्दी संस्करण DOI:
\href{https://doi.org/10.5281/zenodo.21920511}{10.5281/zenodo.21920511}\par}
\vspace{0.35cm}
{\small यह मशीन-सहायित संचयी अनुवाद है। प्रत्येक सम्मिलित खंड का
गणितीय विन्यास जाँचा गया है, PDF सफलतापूर्वक बनाया गया है और प्रत्येक
पृष्ठ का दृश्य परीक्षण किया गया है। यह पूर्ण पुस्तक, मानवीय समीक्षा,
आलोचनात्मक संस्करण अथवा PDF/UA सुगम्यता का दावा नहीं है। अगला अनुवाद
खंड \texttt{trees.tex} है।\par}
\vspace*{0.06\textheight}
\par
\endgroup
\clearpage
\raggedbottom

\tableofcontents*
\clearpage

\subfile{basics}
\subfile{subsets}
\subfile{important-sets}
\subfile{unions-and-intersections}
\subfile{pairs-and-products}
\subfile{russells-paradox}
\subfile{../relations/relations-as-sets}
\subfile{../relations/special-properties}
\subfile{../relations/equivalence-relations}
\begin{sloppypar}
\subfile{../relations/orders}
\end{sloppypar}
\subfile{../relations/graphs}

\end{document}
